\documentclass[11pt]{article}

\usepackage[T1]{fontenc}
\usepackage{lmodern}
\usepackage{amsmath,amssymb,amsthm,mathtools}
\usepackage{booktabs}
\usepackage{array}
\usepackage[margin=1in]{geometry}
\usepackage{microtype}
\usepackage{xcolor}
\usepackage[hidelinks]{hyperref}

\theoremstyle{plain}
\newtheorem{theorem}{Theorem}[section]
\newtheorem{proposition}[theorem]{Proposition}
\newtheorem{corollary}[theorem]{Corollary}
\theoremstyle{remark}
\newtheorem{remark}[theorem]{Remark}

\newcommand{\F}{\mathbb F}
\newcommand{\R}{\mathbb R}
\newcommand{\A}{\mathbb A}
\newcommand{\impl}{\preccurlyeq}
\newcommand{\caff}{C_{\mathrm{aff}}}
\newcommand{\Cone}{\mathcal C}
\newcommand{\Kone}{\mathcal K}
\newcommand{\ra}{\rightarrow}
\newcommand{\source}[1]{{\normalfont\small\textsf{[#1]}}}

\title{Exchange rates as quantitative structure\\on a hierarchy of maps}
\author{Artem Vorozhtsov\\
\small\texttt{artem.voroztsov@gmail.com}}
\date{Working draft, \today}

\hypersetup{pdftitle={Exchange rates as quantitative structure on a hierarchy of maps},
            pdfauthor={Artem Vorozhtsov}}

\begin{document}
\maketitle

\begin{abstract}
Classifying a family of maps under an implementation relation produces a
preorder and stops there.  We add the missing quantitative layer: an
\emph{exchange rate} counting how many copies of one map a second map can
supply per copy, and we study the geometry it induces.  The main result is a
structure theorem --- every exchange geometry is a one-dimensional geometry plus
a defect bounded by $\log 2$, sharply, because the underlying Hilbert projective
metric fibres over a line with fibres of diameter exactly one bit.  We record
what the rate reads off an arithmetic family, why the comparison it induces is
not transitive, and how large that failure is.  Every numbered claim is marked
\emph{proved} or \emph{computed}; sources are the accompanying findings files.
\end{abstract}

\section{The framework is general, and it adds structure rather than objects}

Fix any class of maps $\mathcal P$ and any notion of \textbf{implementation} ---
a relation $f \impl g$ meaning ``$g$ can be wired up to compute $f$''.
Classification of $\mathcal P$ under $\impl$ is the classical activity: it
produces a preorder, and its Hasse diagram is the hierarchy.  For polynomial
maps over a field the standard choice is \emph{affine implementation},
$f = b \circ g \circ a$ with $a,b$ affine, and the resulting posets of quadratic
and cubic maps over $\F_q$ are the objects this project started from.

A preorder answers one question --- \emph{can $g$ do $f$ at all?} --- and is
silent on the quantitative one.  The exchange rate supplies the missing
structure.  Writing $f^{\times k}$ for the $k$-fold Cartesian power and
\[
  k_{g\ra f}(r) \;=\; \max\{\,k : f^{\times k} \impl_{\mathrm{aff}} g^{\times r}\,\},
\]
block-diagonal composition gives $k(r+s) \ge k(r) + k(s)$, so Fekete's lemma
yields
\begin{equation}\label{eq:caff}
  \caff(g\ra f) \;=\; \lim_{r\to\infty} \frac{k_{g\ra f}(r)}{r}.
\end{equation}

\textbf{This is the object.}  It refines the hierarchy without changing it:
$f \impl g$ already implies $\caff(g\ra f) \ge 1$, and $\caff$ grades every
comparable pair by \emph{how many} copies, and every incomparable pair by an
exchange ratio the preorder cannot see.  The construction is generic --- it
needs only a monoidal product and an implementation relation --- so it applies
verbatim to any hierarchy of maps a reader already cares about.

The computable proxy is the \textbf{signature rate}.  The signature $\sigma(f)$
is the multiset of fiber cardinalities; for Cartesian powers fiber sizes
multiply, and the entropy method gives the closed form
\begin{equation}\label{eq:csig}
  C(g\ra f) \;=\; \inf_{0\le\beta\le\infty}
  \frac{\log Z_g(\beta)}{\log Z_f(\beta)},
  \qquad Z_a(\beta) = \sum_i a_i^{\beta}.
\end{equation}
Here $\beta$ is an inverse temperature and $Z$ a partition function with energy
levels $-\log n_i$; \eqref{eq:csig} is a theorem about the operational limit,
not a second definition.  Throughout, $C$ means \eqref{eq:csig} and $\caff$
means \eqref{eq:caff}.

\subsection{The two rates are incomparable}

An earlier draft asserted that affine implementation forces signature
implementation, hence $\caff \le C$.  \textbf{That is false}, and the
counterexample is a one-liner \source{affine\_rate}.  Over $\F_3$ take
$g = x^2+y$ and $f = x^2$.  Then
\[
  f \;=\; g \circ a \qquad\text{with the \emph{singular} affine processor }
  a(x,y) = (x,0),
\]
so $f \impl_{\mathrm{aff}} g$ in one shot and $\caff(g\ra f) = 1$; while
$\sigma(g) = (3,3,3)$ and $\sigma(f) = (6,3)$ give
\begin{equation}\label{eq:incomparable}
  C\bigl(\sigma(g) \ra \sigma(f)\bigr)
  \;=\; \frac{\log 3}{\log 6} \;=\; 0.6131471927654584\ldots \;<\; 1
  \;=\; \caff(g\ra f).
\end{equation}

\textbf{The mechanism.}  The affine model explicitly permits singular input
processors.  Here $a$ is $3$-to-$1$: it collapses the nine-point source onto a
three-point line, and pulling back through it \emph{manufactures} a fiber of
size $6$ from a map all of whose fibers have size $3$.  Signature implementation
cannot do that --- its processors inject each target fiber into an assigned
source fiber, so fiber sizes are inherited and never enlarged.  Neither rate
bounds the other: on the five non-constant affine classes of quadratic maps
$\F_3^2 \ra \F_3$, nine ordered pairs have $\caff < C$ and four have
$\caff > C$.

Worse, $\caff$ is not a function of the signature at all.  The maps $x$ and
$x^2+y$ share the signature $(3,3,3)$, so their signature rates against $x^2$
agree to all $40$ digits; but $\caff(x \ra x^2) = 0$ (the $x$-atoms are affine
functions and span nothing of degree $2$) and $\caff(x^2+y \ra x^2) = 1$.  Two
pairs, identical $C$, affine rates at the two ends of the range.  \emph{No
functor between the two processor categories exists in either direction.}

That the signature loses information was already visible on the arithmetic side:
over $\F_{11}$, $1744$ genus-two pencils $y^2 = P(x)+c$ fall onto only $296$
signatures, and every one of the $209$ multiply-realised signatures comes from
pencils with different fiberwise isogeny data ($419$ of $420$ at $\F_{13}$)
\source{flux\_arithmetic}.  But the affine counterexample is stronger than
lossiness: $C$ is not an approximation to $\caff$ from either side.
\textbf{Every statement below is a statement about $C$, and the word
``signature'' belongs in each of them.}

\section{What the rate reads off an arithmetic family}

Let $f : \A^2 \ra \A^1$ over $\F_q$ have geometrically irreducible fibers,
$N_c = \#f^{-1}(c)$, $a_c = q - N_c$ the Frobenius traces,
$m_2 = q^{-2}\sum_c a_c^2$, and let $L = (q,\dots,q)$ be the flat signature of a
linear map \source{m\_and\_e\_and\_a\_c}.

\begin{proposition}[proved]\label{prop:endpoint}
$C(L\ra f) = \log q / \log \max_c N_c$, attained at $\beta = \infty$.
\end{proposition}

\begin{proposition}[proved]\label{prop:bottleneck}
$1 - C(f\ra L) = (3-2\sqrt2)\,m_2/(2q\log q) + O(q^{-3/2})$, with the infimum at
\[
  \beta^{*} \;=\; \sqrt2 - 1 \;=\; 0.414213562\ldots
\]
independent of family, genus and $q$.
\end{proposition}

So the two ends of the temperature axis read two classical statistics: the
\textbf{extreme Frobenius trace} at $\beta = \infty$ and the \textbf{second
moment} at $\beta^{*}$.  The interior carries the whole moment ladder, each
order damped by ${\approx}\,0.6/\sqrt q$.  Flatness is equivalent to a
permutation-polynomial condition, which subsumes the known $q \equiv 2 \bmod 3$
result and reaches non-congruence conditions such as supersingular primes.

Two things are provably invisible, for structural reasons rather than by
accident: the \emph{smallest} fiber, because isolating it needs $\beta < 0$
where $Z$ is no longer order-preserving; and the Katz--Sarnak symmetry type,
because $\sum_c a_c = 0$ identically and that first moment is exactly the
separating statistic.

\begin{remark}[the objection, stated plainly]
$C(L\ra f)$ is a strictly monotone function of $\max_c N_c$, and $1 - C(f\ra L)$
is to leading order affine in $m_2$.  The rate is a function of the signature
and the signature is the list $\{N_c\}$, so nothing the rate reports about
arithmetic could fail to be computable from $\{N_c\}$ directly.  Information
flows into the matrix, not out of it.  What survives is not a new theorem about
curves but a statement about the \emph{comparison}: of all ways to compress
$\{N_c\}$, resource-theoretic comparison against a linear map returns exactly
the extreme value and the second moment, and returns them at two specific
temperatures.  That is a consistency check on the framework, and it is how it
should be written.
\end{remark}

\section{Main result: one dimension plus a bounded defect}

This is the most important thing the programme has produced, because it is the
only statement about the \emph{framework} rather than about particular
resources.

Write $F_a = \log Z_a$, $u_a = \log F_a$, $R = \log r$ ($r$ = number of fibers),
$\Lambda = \log(\max \text{fiber})$, $\sigma = \log(R/\Lambda)$, and
$s = \log\beta$.

\begin{theorem}[proved; verified to $3.6\cdot10^{-15}$]\label{thm:A}
\source{realizability}
\[
  u_a(e^s) \;=\; \log\Lambda_a + \max(\sigma_a, s) + w_a(s),
  \qquad 0 \le w_a \le \log\!\bigl(1 + e^{-|s-\sigma_a|}\bigr),
\]
with $w_a$ unimodal, peaked exactly at $s = \sigma_a$, and $1$-Lipschitz, of
height $\le \log 2$ with equality iff $a$ is flat.  Every profile is a kink plus
one bump.
\end{theorem}

Splitting $-\log C$ into symmetric and antisymmetric parts --- $S = d/2$ the
exchange metric, $A$ the comparison, with $a \prec b \iff A(a,b) > 0$ ---
Theorem~\ref{thm:A} gives
\begin{align}
  d(a,b) &= |\sigma_a - \sigma_b| + \varepsilon,
  & 0 \le \varepsilon &\le \log\bigl(1+e^{-|\Delta\sigma|}\bigr) \le \log 2,
  \label{eq:metric}\\
  A(a,b) &= \psi(b) - \psi(a) + D,
  & |D| &\le \tfrac12\log\bigl(1+e^{-|\Delta\sigma|}\bigr) \le \tfrac{\log 2}{2},
  \label{eq:comparison}
\end{align}
with $\psi = \frac12\log(\log r \cdot \log M)$.  Both constants are
\textbf{sharp} and both bounds are \textbf{independent of the number of points}.

\begin{theorem}[proved]\label{thm:B}
\source{realizability/OBSTRUCTION}
$d(a,b)$ is exactly the \textbf{Hilbert projective metric} between $F_a$ and
$F_b$ --- of the ambient cone $\Kone$ of positive functions under the pointwise
order, restricted to the achievable set
$\Cone = \{\Phi \text{ convex, nondecreasing}, \Phi \ge \Lambda_\Phi\beta\}$,
which is the projective closure of that set.  Cartesian powers are the
projective rescaling the Hilbert metric quotients out.
\end{theorem}

The distinction between the two cones matters: $d$ is \emph{not} the intrinsic
Hilbert metric of $\Cone$, and reading it that way yields a contraction theorem
that does not exist \source{birkhoff, \S2.5}.

\begin{theorem}[proved]\label{thm:Bprime}
\source{birkhoff}
The constant $\log 2$ is a \textbf{projective diameter}.  With
$\Cone_0 = \{\Phi \in \Cone : \Phi(0) = \Lambda_\Phi\}$ the $\sigma = 0$ fibre,
$\operatorname{diam}_d \Cone_0 = \log 2$ exactly, attained between $\max(1,\beta)$
and $1+\beta$.  The $\beta$-dilation group acts by isometries with geodesic
orbits, $\sigma$ is a surjective $1$-Lipschitz map to $\R$, and every fibre is
isometric to $\Cone_0$.
\end{theorem}

Equivalently: \emph{the exchange geometry is the real line thickened by exactly
one bit.}  That is the geometric source of both sharp constants, and the
attached Birkhoff ratio is $\tanh(\tfrac{\log 2}{4}) = 3 - 2\sqrt2$ exactly.

\subsection{What Theorem \ref{thm:B} does not buy}

No processor is a Birkhoff contraction \emph{(proved)}.  Tensoring is a
\textbf{translation} $\Phi \mapsto \Phi + F_c$, neither linear nor degree-$1$
homogeneous; and every operation that \emph{is} linear shifts $\sigma$ by a
constant, so its projective diameter is infinite and its Birkhoff ratio is
exactly $1$, attained.  Worse, tensoring \textbf{expands} $d$ without bound ---
$d(a, a^{\otimes k}) = 0$ while $d(a\otimes c, a^{\otimes k}\otimes c) \to \log k$
--- so $d$ is a pseudometric on resources and $\otimes$ does not descend to its
quotient.  Any claim that conversion loses distinguishability at a definite rate
is false.  The correct positive statement is that all four resource operations
are \textbf{nonexpansive in Thompson's part metric}: Hilbert and Thompson differ
exactly along the projective direction, and the projective direction is what
conversion trades in.

\subsection{Everything previously observed is a corollary}

The $\ell_2$-distortion ${\approx}\,1.1$ with $c_2/\log n$ \emph{falling} ---
Bourgain's $O(\log n)$ is unreachable because the non-line budget does not grow
with $n$.  The mild pentagonal failure, and every negative-type violation,
confined to a window of diameter $2\log 2$.  The measured $4$--$9\%$ curl of the
comparison and its flatness in $n$ from $8$ to $698$ points --- the bounds in
\eqref{eq:comparison} are per-edge.  The same constant reappearing in the
$q\to\infty$ limit.

\begin{corollary}[proved]\label{cor:C}
The geometric mean of the rate asymmetry $C(b\ra a)/C(a\ra b)$ around
\emph{any} preference cycle is at most $2$.
\end{corollary}

Sharpness is closed on one side.  Writing $\operatorname{curl} A = \frac12\log\Omega$
with $\Omega$ the ratio of forward to backward arbitrage around the triangle
\emph{(proved)}, the question ``is the sharp constant $(\log 2)/2$?'' is exactly
``\emph{is $\Omega \le 2$} --- does the forward arbitrage exceed the backward by
at most one bit?''  The lower bound is proved in closed form by
$\Phi_1 = \max(1,\beta)$, $\Phi_2 = 1+\beta$, $\Phi_3 = 1+e^T\beta$, giving
$\operatorname{curl} A = \frac12(\log 2 - \log(1+e^{-T})) \nearrow (\log 2)/2$ ---
so the supremum is $\ge (\log 2)/2$ and is \emph{not attained}
\source{birkhoff}.  The upper bound remains open between $(\log 2)/2$ and
$3(\log 2)/2$, and any proof must use the $\sigma$-spread: inside a single fibre
the supremum is only $\frac12\log(4/3)$.

\section{Are the cycles important?}

The exchange comparison is not transitive.  Over $\F_{11}$ there are $132$
strict $3$-cycles among genus-two pencils $y^2 = P(x)+c$ in a complete
enumeration, $1475$ over $\F_{13}$, with certified witnesses over $\F_{11}$ and
$\F_{101}$ verified three ways.  Substantial effort went into finding and
certifying them.  \textbf{Are they worth it?}  Three results weaken the obvious
reading, one strengthens it, and one bounds how much can be claimed.

\paragraph{Against: cycles are not arithmetic.}
A random pool of signatures matched to a curve pool in the only things that
matter ($q$ fibers, $\sum_c N_c = q^2$, matched trace distribution) has the
\emph{same} curl, and the tightest control has slightly more.  Over $108$ pools,
$\lVert\operatorname{curl}\rVert/\lVert A\rVert$ is a function of the trace
spread to $R^2 = 0.93$ \source{flux\_arithmetic}.  Cycles are a property of
near-flat signature \emph{shape}, not of the arithmetic that produced it.

\paragraph{Against: cycles mostly do not descend to the maps.}
The $\F_3$ tensor three-cycles are cycles of \emph{signatures}, and since
$\caff$ is not a signature invariant (\S1) the question is only well posed after
choosing affine orbits.  The seven signature cycles lift to twelve orbit
triples, and \textbf{eight of the twelve are refuted under $\caff$}
\source{affine\_rate} --- in each case because one vertex is affinely
implementable from another, forcing the opposite comparison on that edge.  Where
the descent can be tested, it mostly fails.

\paragraph{For, and this replaces an argument we had wrong.}
An earlier draft claimed cycles are a zero-temperature artefact --- that
averaging over temperature destroys them.  That is false \source{maslov}.  Write
$c_\beta(a\ra b) = F_a(\beta)/F_b(\beta)$ for the fixed-$\beta$ exchange rate and
let $M_\lambda$ be the $\lambda$-power mean over a temperature prior $\rho$.
Then the whole comparison is
\begin{equation}\label{eq:powermean}
  A_\lambda(a,b) \;=\; \tfrac12\log\frac{C_\lambda(a\ra b)}{C_\lambda(b\ra a)},
  \qquad C_\lambda(a\ra b) = M_{-\lambda}\bigl(c_{\cdot}(a\ra b)\bigr),
\end{equation}
with $\lambda = \infty$ the framework's infimum, $\lambda = 1$ the plain
arithmetic average of the transition rates, and $\lambda = 0$ the geometric
mean.  Two facts:
\begin{itemize}
\item At a \emph{single} temperature the comparison is the total order of the
  scalar $F_{\cdot}(\beta)$, since $c_\beta(a\ra b)c_\beta(b\ra a) = 1$.  No pool
  cycles at any one temperature; all non-scalar content comes from the
  $\beta$-dependence of the argmin.
\item Cycles survive averaging.  There are certified $3$-cycles of \emph{flat}
  signatures --- which Theorem~\ref{thm:A} \emph{proves} cannot cycle at
  $\lambda = \infty$ --- living in bounded bands
  $\lambda \in (1.5211, 2.5376)$ and $(0.8032, 1.3139)$, the second a cycle at
  $\lambda = 1$ exactly.
\end{itemize}
The comparison is a potential \emph{iff} $\lambda = 0$: the geometric mean is
the unique power mean commuting with reciprocals, so
$C_0(a\ra b)C_0(b\ra a) = 1$ and the trade becomes reversible.  But the same
identity gives $S_\lambda = \lambda\operatorname{Var}_\rho/2 + O(\lambda^3)$, so
at $\lambda = 0$ every distance is zero and every pair freely interconvertible.
The scalar complexity available there is the potential of an empty theory.

\begin{remark}
Non-scalarity is not an artefact of the infimum.  It is a consequence of
\textbf{irreversibility}: the comparison fails to be a potential exactly because
a power mean of order $\lambda \ne 0$ does not commute with reciprocals.  The
infimum is the extreme case, not the source.
\end{remark}

\paragraph{Where the cycles come from, mechanically.}
Brand\~ao--Horodecki--Ng--Oppenheim--Wehner govern conversion under thermal
operations by a continuum of R\'enyi $\alpha$-free energies --- the same
$\beta$-indexed family.  Their scalar is an infimum of a \textbf{difference},
hence a partial order, acyclic by construction; ours is an infimum of a
\textbf{ratio}, hence a tournament, which can cycle \source{quantum}.  And the
ratio is \emph{forced}: these resources are unnormalised, so the difference
diverges like $\beta(\Lambda_a - \Lambda_b)$ and the ratio is the only finite
scalar.  \textbf{Cycles are the price of comparing resources with no common
normalisation} --- the situation for maps, where nothing plays the role of a
fixed particle number.

\paragraph{The hedge, and it is the number to quote.}
A single \emph{unfitted} temperature prior --- uniform in $s = \log\beta$ on
$[0,6]$ --- reproduces $99.39\%$ of the tropical order on the $\F_{11}$
arithmetic pool and $99.51\%$ on $\F_{13}$; a fitted prior reaches $99.62\%$ and
transfers across $q$ without loss, against an information-theoretic ceiling of
$99.90\%$ (at least $42$ edge-disjoint $3$-cycles must be misordered).  Both
beat the HodgeRank potential $\psi_{\mathrm{opt}}$ at $98.31\%$
\source{maslov}.  So the irreducibly non-scalar part of the comparison is about
\textbf{four parts in a thousand}.

\paragraph{The defensible claim.}
Asymptotic convertibility is governed by no scalar complexity, because the
comparison is a non-reversible aggregation over temperature.  The failure is
real but small: a single temperature prior recovers $99.4\%$ of the order, and
at $\lambda = \infty$ rates around any preference cycle differ by a
geometric-mean factor of at most two.  What this does \emph{not} support is the
story the cycle hunt was implicitly aimed at --- that cycles reveal something
arithmetic.  They do not.

\paragraph{One genuinely quantum instance.}
Replacing $\operatorname{Tr}A^\beta$ by a sandwiched R\'enyi divergence against
a background operator gives a non-spectral rate under which three qubit states
form a certified $3$-cycle
($|\operatorname{curl}A|/\sum|A| = 1$ exactly, $40$ digits) whose decohered
shadow is transitive \source{quantum}.  Two caveats: the plain quantisation
$\operatorname{Tr}A^\beta$ is \emph{empty} --- it factors through the spectrum
identically, so a quantum cycle without a classical shadow cannot exist --- and
every sandwiched profile still lies in the classical cone.  What is certified is
the narrower physical statement that decoherence destroys that cycle, not that
the geometry is new.

\section{Status}

Proved and independently re-verified: Propositions~\ref{prop:endpoint} and
\ref{prop:bottleneck}, Theorems~\ref{thm:A}--\ref{thm:Bprime},
Corollary~\ref{cor:C}, the power-mean identity \eqref{eq:powermean} and its
$\lambda = 0$ reversibility, the flatness/permutation-polynomial theorem, the
invisibility of the smallest fiber and of symmetry type, superadditivity of $C$
and $\caff$ under Cartesian products --- and, for $\caff$, \textbf{strict}
superadditivity, witnessed by $x^2 \otimes x \ra x^2+y$ over $\F_3$ with a gap of
at least $1/2$ \source{affine\_rate}.  That settles the one form of the synergy
question no variational formula could have decided: one affine processor draws a
degree-$2$ part from $x^2$ and a linear part from $x$, which neither factor
supplies alone.

The first exact affine rates are in hand --- $1/2$ for $x^2\ra x$, $x^2\ra xy$,
$x^2\ra x^2+y^2$, and $2/3$ for $x^2+y^2\ra x$ --- and in the last the Fekete
supremum is attained at no $r \le 2$, so no single-shot computation determines
$\caff$.

Computed and certified: the cycle censuses and their witnesses; the
flat-signature band-cycles; the qubit cycle.

\paragraph{Open and load-bearing.}
Prove $\Omega \le 2$, equivalently $|\operatorname{curl}A| \le (\log 2)/2$;
whether $\caff$ cycles on affine orbits; whether any $\lambda < \infty$ is
operational; six of the twenty-five affine brackets over $\F_3$.

\paragraph{Withdrawn.}
Three claims of earlier drafts.  That the exchange metrics do not fill
$\mathrm{MET}$ --- $C_4$ is realisable exactly, and the obstruction rested on a
stalled search.  That $\caff \le C$ --- false, with the singular-processor
counterexample in \eqref{eq:incomparable}.  And that cycles are a
zero-temperature artefact --- false, with certified cycles at $\lambda = 1$.
Each was refuted by the session sent to verify it.

\end{document}
